\section{Tool Support}
\subsection{One checker, three input routes}
The command-line tool accepts symbolic text and versioned Core JSON. Both reach
the same checker; reports distinguish conformance failures, binding/decoding
errors and unsupported inputs. Figure~\ref{fig:boundary} locates the proved
connections. Equation~\ref{eq:adequacy} begins at the AST, outside the parser.
The JSON decoder admits malformed represented records for the checker to reject;
its key-normalization policy is documented separately.

\begin{figure}[t]
\centering
\begin{tikzpicture}[font=\scriptsize,
 box/.style={draw,align=center,minimum width=20mm,minimum height=8mm},
 arr/.style={-{Stealth[length=1.5mm]}}]
\node[box] (text) at (-4,1) {source text};
\node[box] (ast) at (-1.3,1) {source AST};
\node[box] (json) at (-4,0) {Core JSON};
\node[box] (decode) at (-1.3,0) {decoder};
\node[box] (native) at (-4,-1) {Ecore/XMI};
\node[box] (map) at (-1.3,-1) {load and map};
\node[box,minimum height=29mm,fill=black!4] (core) at (1.5,0) {Core schema\\and snapshot};
\node[box] (check) at (4.2,0) {decision};
\draw[arr] (text) -- node[above]{parse} (ast);
\draw[arr,double,double distance=.5pt] (ast) -- node[above]{elaborate} ([yshift=10mm]core.west);
\draw[arr] (json) -- (decode);
\draw[arr] (decode) -- (core);
\draw[arr] (native) -- (map);
\draw[arr] (map) -- ([yshift=-10mm]core.west);
\draw[arr,double,double distance=.5pt] (core) -- node[above]{check} (check);
\node[box] (export) at (-4,-2.7) {Core export};
\node[box] (fresh) at (-1.3,-2.7) {fresh Ecore/XMI};
\node[box] (reload) at (1.5,-2.7) {EMF reload};
\node[box] (compare) at (4.2,-2.7) {mapped comparison};
\draw[arr] (export) -- (fresh);
\draw[arr] (fresh) -- (reload);
\draw[arr] (reload) -- (compare);
\node[align=center] at (.1,-3.4) {Native experiment: compare original and reloaded observations\\under the retained identity bijection.};
\end{tikzpicture}
\caption{Input routes and the separate native round-trip experiment. Double
arrows identify computations covered by formal correspondence results under
their stated premises. Parsing, codecs, loading and native comparison are tested
boundaries. An actual EMF validator is a separate experimental oracle.}
\label{fig:boundary}
\end{figure}

\subsection{Restricted EMF realization}
The Java bridge loads Ecore/XMI through dynamic EMF objects. Its restricted
profile rejects unsupported annotations, operations, generics, declared defaults,
derived/unsettable features, custom datatypes and unresolved references. Export
also requires class-owned association ends, representable packages/names and
signed 32-bit integers.

Import allocates identities from containment. An XML pass tracks written scalar
values separately from omitted ones. This lexical policy can differ from native
\texttt{eIsSet}, especially for enumeration defaults; Section~5 reports an actual
mismatch. Lossless alignment is therefore checked per input, not assumed from
successful import.

Export constructs fresh resources, establishes both opposite pointers and compares
observations before serialization. A sidecar maps Core identities to native URIs.
Reload comparison requires a bijection, remaps references, and compares sequences
or multisets according to ordering. These are executable checks outside the Lean
proof, not universal codec correctness results.

\subsection{Reuse for bounded metadata}
A fixed metadata vocabulary represents class and scalar-attribute descriptions
as ordinary objects and properties. The common checker validates a description
before an interpreter produces a Core schema. The example describes an optional
Boolean property and then checks an instance of the resulting class. A dangling
owner is rejected as metadata; reversed bounds can be well-typed metadata yet
produce an invalid interpreted schema. This separates validity of the
description graph from validity of what it denotes.
